\documentclass[12pt]{article}
\addtolength{\topmargin}{-2cm}
\addtolength{\textheight}{3cm}
\addtolength{\oddsidemargin}{-1.5cm}
\addtolength{\textwidth}{2cm}
\usepackage{graphicx}
\pagestyle{empty}
\begin{document}
\thispagestyle{empty}
\hspace*{-20pt}
\includegraphics[scale=0.8,angle=0]{iiser_logo.eps}
\begin{minipage}[b]{13.5cm}
\begin{center}
\vspace*{-11pt}
\baselineskip=18pt
{\Large \bf Indian Institute of Science Education\\
\& Research Mohali}\\
{\bf Sector 81, Mohali, Phone \& Fax:+91-172-2790188 
Email:~arvind@iisermohali.co.in}
\vspace*{6pt}
\end{center}
\end{minipage}
\hrule
\vspace*{10pt}
\rightline{Prof. Arvind}
\rightline{May 7, 2018}
\vspace*{10pt}
\begin{flushleft}
Professor Matteo Paris, \\
Editor,\\
Physics Letters A\\
\end{flushleft}
\vspace*{6pt}
\noindent
{\bf Subject:} Manuscript submission\\

\vspace*{6pt}
\noindent
{Dear Professor Paris,}

\vspace*{6pt}
  
We are herewith submitting a paper titled ``Revisiting the
admissibility of non-contextual hidden variable models in
quantum mechanics'' authored by Atul Singh Arora, Kishor
Bharti and Arvind.


In this work we propose a new (to the best of our knowledge)
non-contextual hidden variable model which helps clarify the
limitations of a well-established notion in the foundations
of quantum mechanics, that of contextuality. Contextual
hidden variable models on the other hand are already known
but these are in stark contrast with our model.
Understanding the notion of contextuality and its
implications better is therefore likely to stimulate further
research. Our work contributes towards this older direction
which has witnessed a renewed interest.

Thanking you, 


\vspace*{6pt}
\begin{flushright}
Yours Sincerely, \\
\vspace*{12pt}
Arvind\\
\end{flushright}
\end{document}

